\begin{lstlisting}[language=c++]
  #include "BPatch.h"
  #include "BPatch_addressSpace.h"
  #include "BPatch_binaryEdit.h"
  #include "BPatch_function.h"
  #include "BPatch_point.h"
  #include "BPatch_process.h"

  #include <cstlib>
  #include <iostream>
  #include <vector>

  // Create an instance of class BPatch
  BPatch bpatch;

  // Different ways to perform instrumentation
  enum accessType_t { create, attach, open };

  // Attach, create, or open a file for rewriting
  BPatch_addressSpace *startInstrumenting(accessType_t accessType,
                                          const char *name, int pid,
                                          const char *argv[]) {
    switch (accessType) {
    case create: {
      auto handle = bpatch.processCreate(name, argv);
      if (!handle) {
        fprintf(stderr, "processCreate failed\n");
      }
      return handle;
    }
    case attach: {
      auto handle = bpatch.processAttach(name, pid);
      if (!handle) {
        fprintf(stderr, "processAttach failed\n");
      }
      return handle;
    }
    case open:
      // Open the binary file and all dependencies
      auto handle = bpatch.openBinary(name, true);
      if (!handle) {
        fprintf(stderr, "openBinary failed\n");
      }
      return handle;
    }
    return nullptr;
  }

  // Find a point at which to insert instrumentation
  std::vector<BPatch_point *> *findPoint(BPatch_addressSpace *app,
                                         const char *name,
                                         BPatch_procedureLocation loc) {
    BPatch_image *appImage = app->getImage();

    // Scan for functions named "name"
    std::vector<BPatch_function *> functions;
    appImage->findFunction(name, functions);
    if (functions.size() == 0) {
      fprintf(stderr, "No function %s\n", name);
      return nullptr;
    }

    if (functions.size() > 1) {
      fprintf(stderr, "More than one %'s;' using the first one\n", name);
    }

    // Locate the relevant points
    return functions[0]->findPoint(loc);
  }

  // Create and insert an increment snippet
  bool createAndInsertSnippet(BPatch_addressSpace *app,
                              std::vector<BPatch_point *> *points) {
    BPatch_image *appImage = app->getImage();

    // Create an increment snippet
    BPatch_type *type = appImage->findType("int");
    BPatch_variableExpr *counter = app->malloc(*type, "myCounter");
    BPatch_arithExpr addOne(
        BPatch_assign, *counter,
        BPatch_arithExpr(BPatch_plus, *counter, BPatch_constExpr(1)));

    // Insert the snippet
    if (!app->insertSnippet(addOne, *points)) {
      fprintf(stderr, "insertSnippet failed\n");
      return false;
    }
    return true;
  }
  // Create and insert a printf snippet
  bool createAndInsertSnippet2(BPatch_addressSpace *app,
                               std::vector<BPatch_point *> *points) {
    BPatch_image *appImage = app->getImage();

    // Create the printf function call snippet
    std::vector<BPatch_snippet *> printfArgs;
    BPatch_snippet *fmt =
        new BPatch_constExpr("InterestingProcedure called %d times\n");
    printfArgs.push_back(fmt);
    BPatch_variableExpr *var = appImage->findVariable("myCounter");
    if (!var) {
      fprintf(stderr, "Could not find 'myCounter' variable\n");
      return false;
    }
    printfArgs.push_back(var);

    // Find the printf function
    std::vector<BPatch_function *> printfFuncs;
    appImage->findFunction("printf", printfFuncs);
    if (printfFuncs.size() == 0) {
      fprintf(stderr, "Could not find printf\n");
      return false;
    }

    // Construct a function call snippet
    BPatch_funcCallExpr printfCall(*(printfFuncs[0]), printfArgs);

    // Insert the snippet
    if (!app->insertSnippet(printfCall, *points)) {
      fprintf(stderr, "insertSnippet failed\n");
      return false;
    }
    return true;
  }
  void finishInstrumenting(BPatch_addressSpace *app, const char *newName) {
    BPatch_process *appProc = dynamic_cast<BPatch_process *>(app);
    BPatch_binaryEdit *appBin = dynamic_cast<BPatch_binaryEdit *>(app);
    if (appProc) {
      if (!appProc->continueExecution()) {
        fprintf(stderr, "continueExecution failed\n");
      }
      while (!appProc->isTerminated()) {
        bpatch.waitForStatusChange();
      }
    } else if (appBin) {
      if (!appBin->writeFile(newName)) {
        fprintf(stderr, "writeFile failed\n");
      }
    }
  }

  int main() {
    // Set up information about the program to be instrumented
    const char *progName = "InterestingProgram";
    int progPID = 42;
    const char *progArgv[] = {"InterestingProgram", "- h", NULL};
    accessType_t mode = create;

    // Create/attach/open a binary
    BPatch_addressSpace *app =
        startInstrumenting(mode, progName, progPID, progArgv);
    if (!app) {
      fprintf(stderr, "startInstrumenting failed\n");
      return EXIT_FAILURE;
    }
    // Find the entry point for function InterestingProcedure
    const char *interestingFuncName = "InterestingProcedure";
    std::vector<BPatch_point *> *entryPoint =
        findPoint(app, interestingFuncName, BPatch_entry);
    if (!entryPoint || entryPoint->size() == 0) {
      fprintf(stderr, "No entry points for %s\n", interestingFuncName);
      return EXIT_FAILURE;
    }
    // Create and insert instrumentation snippet
    if (!createAndInsertSnippet(app, entryPoint)) {
      fprintf(stderr, "createAndInsertSnippet failed\n");
      return EXIT_FAILURE;
    }
    // Find the exit point of main
    std::vector<BPatch_point *> *exitPoint = findPoint(app, "main", BPatch_exit);
    if (!exitPoint || exitPoint->size() == 0) {
      fprintf(stderr, "No exit points for main\n");
      return EXIT_FAILURE;
    }
    // Create and insert instrumentation snippet 2
    if (!createAndInsertSnippet2(app, exitPoint)) {
      fprintf(stderr, "createAndInsertSnippet2 failed\n");
      return EXIT_FAILURE;
    }
    // Finish instrumentation
    const char *progName2 = "InterestingProgram - rewritten ";
    finishInstrumenting(app, progName2);
  }

\end{lstlisting}
